\section{Data and Experimental Setup}
\label{sec:data}

This section records how each panel was assembled and how the method was configured, in enough detail to reproduce the runs. A single principle governs the target across all five panels. Every entity is placed on a shared time grid, and an entity is admitted only when it is observed over the whole window, or over almost all of it, so that the trajectory being clustered is measured rather than imputed. Where a small fraction of interior points is missing, those points are filled by linear interpolation between observed values, and entities missing more than a set fraction of the window are dropped rather than reconstructed. No entity is removed by random subsampling in any run, so each reported segmentation is fitted on the full admitted set.

\paragraph{Carbon dioxide emissions.} The source is the Our World in Data emissions dataset \citep{owidco2}. An entity is a country identified by its ISO code, and aggregate rows that the source marks with non-country codes are excluded so that only sovereign countries enter. The target is annual carbon dioxide emissions per capita in tonnes. The window runs yearly from 1990 to 2022, which is thirty-three time points, and a country is admitted only when its emission series is complete over every year of that window, which yields 208 countries. The covariates, each supplied to the temporal feature construction, are gross domestic product per capita, population, and the shares of coal, oil, and gas in the energy mix.

\paragraph{Pandemic mortality.} The source is the Our World in Data COVID-19 dataset \citep{owidcovid}. An entity is a country, and the target is weekly deaths per million, obtained by aggregating the daily smoothed series to calendar weeks. The time axis is aligned to the epidemic rather than to the calendar. Week zero for a country is the first week in which its cumulative deaths per million exceed one, so that the trajectories are compared from a common epidemiological start rather than from a shared date, and the window is the 104 weeks, roughly two years, that follow that onset. A country is admitted only when it is observed for the full 104-week horizon, which yields 224 countries. The covariates are national attributes that are fixed over the window, namely median age, population, population density, gross domestic product per capita, diabetes prevalence, hospital beds per thousand, and life expectancy.

\paragraph{National income.} The source is the World Bank World Development Indicators \citep{worldbankwdi}, retrieved through the institution's data interface. An entity is a country, and the target is gross domestic product per capita expressed relative to its 1990 value, so that each trajectory begins at one and records relative growth rather than absolute level. The window runs yearly from 1990 to 2023, which is thirty-four time points, and a country is admitted only when its series is complete over every year, which yields 190 countries. The covariates are birth rate, death rate, access to electricity, internet use, and trade in services.

\paragraph{Fine particulate air quality.} The source is the United States Environmental Protection Agency air-quality monitoring archive \citep{epaairdata}. An entity is a monitoring site, and the target is the monthly average concentration of fine particulate matter. The series are placed on a monthly grid spanning 2010 to 2024, which is 180 time points, and a site is admitted when at least ninety percent of the months are observed, with the remaining interior gaps interpolated, which yields 526 sites. The covariates are the latitude, longitude, and elevation of the site together with categorical descriptors of its land use and its setting, such as residential, industrial, or rural, and these categorical descriptors are one-hot encoded so that a rule can test a single category.

\paragraph{Household health spending.} The source is the Consumer Pyramids Household Survey collected by the Centre for Monitoring Indian Economy. An entity is a household, and the target of a household in a month is the logarithm of one plus the ratio of its health spending to its income. The construction of this target and the window are set out in Section~\ref{sec:cmie}, since they involve choices specific to a ratio target, namely a floor on the income denominator, the removal of data-collection sentinels, a fixed ninety-month window centred on the onset of the pandemic, and an admission rule that keeps households observed over at least ninety percent of that window. The full admitted cohort of 14652 households is used, and reporting is confined to subgroups and rules so that no household-level record is disclosed. The surveyed households are not a population sample, so the findings describe the sample rather than the country.

\paragraph{Method configuration.} All five runs use one configuration so that the studies are comparable. The randomised search draws five thousand candidate groupings. The number of groups is drawn between two and eight, the decision tree depth between one and four, and the minimum leaf size between two and thirty percent of the entities. The complexity weight $\lambda$ in Equation~\ref{eq:cost} is set to $0.005$, and the diverse shortlist retains a candidate only when its adjusted Rand similarity to every candidate already retained stays below $0.8$. The trajectory distance is computed on cumulative trajectories as described in Section~\ref{sec:distance}, and the descriptor matrix combines static attributes with the engineered temporal features of Section~\ref{sec:temporal}.
